% Part: first-order-logic 
% Chapter: axiomatic-deduction 
% Section: rules-and-proofs

\documentclass[../../../include/open-logic-section]{subfiles}

\begin{document}

\iftag{FOL}
      {\olfileid{fol}{axd}{rul}}
      {\olfileid{pl}{axd}{rul}}

\olsection{Regels en \usetoken{P}{derivation}}

\begin{explain}
  Axiomatische !!{derivation}s horen bij de eenvoudigste stelsels van
  afleidingsregels in de logica. !!^a{derivation} is gewoon een rij
  !!{formula}s. Elke !!{formula} in die rij moet aan een van twee
  voorwaarden voldoen. De formule is een invulling van een axioma, dus
  van een vaste grondaanname van het formele systeem. Of de formule
  volgt volgens een afleidingsregel uit een of meer eerdere
  !!{formula}s in de rij. De laatste !!{formula} in de rij is wat de
  !!{derivation} !!{derive}s.
\end{explain}

\begin{defn}[!!^{derivability}]
Neem een verzameling~$\Gamma$ van !!{formula}s van $\Lang L$. Een
\emph{!!{derivation}} uit~$\Gamma$ is een eindige rij $!A_1$,
\dots,~$!A_n$ van !!{formula}s. Bij elk nummer $i \le n$ moet een van
de volgende gevallen gelden:
\begin{enumerate}
\item $!A_i \in \Gamma$.
\item $!A_i$ is een axioma.
\item $!A_i$ volgt volgens een afleidingsregel uit een formule $!A_j$
  (en soms ook $!A_k$). Daarbij geldt $j < i$ (en $k < i$): de
  benodigde formules staan dus eerder in de rij.
\end{enumerate}
\end{defn}

Of een !!{derivation} correct is, hangt af van de regels die we
gebruiken en van wat we als axioma nemen. Een afleidingsregel heeft de
vorm ``als dit geldt, dan geldt dat.'' Hij zegt dat stap~$A_i$ in
!!a{derivation} correct is zodra aan bepaalde voorwaarden is voldaan.

\begin{defn}[Afleidingsregel]
Een \emph{afleidingsregel} geeft een voorwaarde die garandeert dat een
stap in !!a{derivation} uit~$\Gamma$ correct is.
\end{defn}

Begin met het eenvoudigste geval. De rij die alleen uit $!A$ bestaat,
met $!A \in \Gamma$, is zonder meer !!a{derivation}. We kunnen dus de
volgende eenvoudige afleidingsregel gebruiken:
\begin{quote}
  Als $!A \in \Gamma$, dan is $!A$ altijd een correcte stap in elke
  !!{derivation} uit~$\Gamma$.
\end{quote}
Hetzelfde geldt als $!A$ een van de axioma's is. Dan is de rij die
alleen uit $!A$ bestaat ook !!a{derivation}. Dat geeft deze regel:
\begin{quote}
  Als $!A$ een axioma is, dan is $!A$ een correcte stap.
\end{quote}
Interessanter wordt het als een regel terugkijkt naar !!{formula}s die
al eerder in de rij staan. De volgende regel heet \emph{modus ponens}:
\begin{quote}
  Als $!B \lif !A$ en $!B$ eerder in de !!{derivation} staan, dan
  telt~$!A$ als een correcte stap.
\end{quote}
Stel dat dit de enige afleidingsregel is. Dan zegt onze definitie
hierboven precies het volgende: $!A_1$, \dots,~$!A_n$ is
!!a{derivation} als en alleen als bij elke $i \le n$ een van deze
gevallen geldt:
\begin{enumerate}
\item $!A_i \in \Gamma$.
\item $!A_i$ is een axioma.
\item voor een zekere $j < i$ is $!A_j$ de formule $!B \lif !A_i$, en
  voor een zekere $k < i$ is $!A_k$ de formule~$!B$.
\end{enumerate}
Het laatste geval zegt dat $!A_i$ met modus ponens volgt uit~$!A_j$
($!B \lif !A_i$) en $!A_k$ ($!B$). We kunnen de hele rij dus stap voor
stap controleren, van $1$ tot~$n$. Bij elk nummer kijken we of de
!!a{formula}~$!A_i$ in~$\Gamma$ staat, een axioma is of volgens een
afleidingsregel een correcte stap is. Als dat bij elk nummer lukt, telt
de hele rij als een correcte !!{derivation}.

\begin{defn}[!!^{derivability}]
We zeggen dat !!a{formula}~$!A$ \emph{!!{derivable}} is uit $\Gamma$.
We schrijven dit als $\Gamma \Proves !A$. Dat geldt als we uit
$\Gamma$ !!a{derivation} kunnen maken waarvan de laatste formule $!A$
is.
\end{defn}

\begin{defn}[Stellingen]
!!^a{formula}~$!A$ heet een \emph{stelling}. Daarvoor moet er uit de
lege verzameling !!a{derivation} bestaan die eindigt in~$!A$. We
schrijven $\Proves !A$ als $!A$ een stelling is en $\Proves/ !A$ als
dat niet zo is.
\end{defn}


\end{document}
